\documentclass[10pt,a4paper]{amsart}
\usepackage[margin=19mm]{geometry}
\usepackage{amsmath,amssymb,mathtools}
\usepackage[scr=rsfs,cal=euler]{mathalfa}
\usepackage{xfrac}
\usepackage[unicode,hidelinks,bookmarks=false]{hyperref}
\newcommand{\ie}{i.e.\ }
\newcommand{\cf}{cf.\ }
\newcommand{\resp}{respectively\ }
\newcommand{\Subsection}{\subsection}
\providecommand{\nobreakdash}{\nobreak\hbox{-}}
\setlength{\emergencystretch}{2em}
\multlinegap0pt
\usepackage{bm}
\usepackage{bbm}
\usepackage{dsfont}
\usepackage{xspace}
\usepackage{cancel}
\usepackage{mathtools}
\usepackage{amsthm}
\makeatletter
\@ifundefined{corollary}{\newtheorem{corollary}{Corollary}}{}
\makeatother
\newcommand{\DI}{\tilde{X}}	% A manifold
\newcommand{\MA}{X}	% The state space
\newcommand{\RR}{\mathbb{R}}	% The reals
\newcommand{\bscdot}{\boldsymbol{\cdot}}
\newcommand{\di}{x}	% An element in DI
\newcommand{\lie}{\mathcal{L}}	% Lie derivative
\title[Sliding motions on systems with non-Euclidean state spaces: ]{Sliding motions on systems with non-Euclidean state spaces: A differential-geometric perspective}
\author{Fernando Casta\~nos}
\date{Source: arXiv:2504.19504v2 (CC BY 4.0)}
\begin{document}
\maketitle

\noindent\emph{Source and licence.} Sliding motions on systems with non-Euclidean state spaces: A differential-geometric perspective. Version arXiv:2504.19504v2 (CC BY 4.0), distributed under the Creative Commons Attribution 4.0 International licence, as recorded in the arXiv licence metadata for this exact version and shown on the abstract page https://arxiv.org/abs/2504.19504v2. Full rights notice: licenses/arxiv-2504.19504-CC-BY-4.0.txt.\par\smallskip

\noindent\textit{Editorial scope.} Counted unit: the Corollary whose source label is \texttt{cor:descend\_vf} in the article (source0.tex lines 344--352, source label \texttt{cor:descend\_vf}), delivered together with the article's own complete original proof of it (source0.tex lines 354--370, one \texttt{proof} environment of 17 lines). No mathematics is written, repaired or re-derived: statement and proof are the article's own text line for line. Required context retained uncounted: eq:descend (lines 284--286); eq:descend\_vf\_1 (lines 337--339). Every definition, display and label of those lines is retained unchanged. Omitted from this package and not redistributed: 1--283, 287--336, 340--343, 353--353, 371--1119. Nothing inside a retained excerpt is omitted or altered: every retained body is delivered exactly as the article's lines are written, so no omission receipt of any kind accompanies this package. Citations: the retained excerpts carry no citation command at all, so no bibliography entry of the article is transcribed and no reference is redistributed. Reference adaptation: none was needed and none was made; every reference inside the delivered unit resolves inside it, so no unresolved reference or citation is produced. Printed slips kept literal: no printed word, symbol, hypothesis or number of the counted statement or of its proof was altered. Layout and class: the article's own class is \texttt{article} with packages including graphicx, tikz-cd, amsmath, amsthm, amsfonts, amssymb, enumerate, color; the wrapper is the lane's pinned \texttt{amsart} editorial wrapper at 10pt on A4, which restates the theorem-like environments the retained text needs and adds the article's own macro definitions that the retained text needs (\texttt{\textbackslash DI}, \texttt{\textbackslash MA}, \texttt{\textbackslash RR}, \texttt{\textbackslash bscdot}, \texttt{\textbackslash di}, \texttt{\textbackslash lie}), with extra wrapper packages (bm, bbm, dsfont, xspace, cancel, mathtools). The delivered numbering is the unit's own. Assets: no retained span contains a figure environment, a graphics-inclusion command, a PDF-page inclusion, a table or a TikZ picture; no image file is copied into this package, the package declares no asset dependency and no image file is redistributed with it. Licence: version arXiv:2504.19504v2 of this article is distributed under the Creative Commons Attribution 4.0 International licence, as recorded in the arXiv licence metadata for this exact version and shown on the abstract page https://arxiv.org/abs/2504.19504v2. A full rights notice is delivered at licenses/arxiv-2504.19504-CC-BY-4.0.txt. Characters: every retained excerpt is pure ASCII, so the delivered engine, XeLaTeX with its default Latin Modern text font, reports no missing character.
 \begin{equation} \label{eq:descend}
  \tilde{s}(\di) = \tilde{s}(\gamma\bscdot\di) \quad \text{for all $\di \in \DI$ and $\gamma \in \Gamma$} \;,
 \end{equation}

\begin{equation} \label{eq:descend_vf_1}
 q_* \tilde{f}(\di,t) = f(q(\di),t) \quad \text{for all $(\di,t) \in \DI\times\RR$} \;, 
\end{equation}

\begin{corollary}[Vector Field on Quotient Manifold] \label{cor:descend_vf}
 If $\tilde{f}$ is a smooth time-dependent vector field on $\DI$ such that, for any smooth real-valued function $\tilde{s}$
 satisfying~\eqref{eq:descend}, we have that
 \begin{equation} \label{eq:descend_vf}
  \lie_{\tilde{f}}\tilde{s}(\di) = 
   \lie_{\tilde{f}}\tilde{s}(\gamma\bscdot\di) \quad \text{for all $(\di,t) \in \DI\times\RR$ and $\gamma \in \Gamma$} \;,
 \end{equation}
 then there exists a unique smooth time-dependent vector field $f$ on $\MA$ that is $q$-related to $\tilde{f}$.
\end{corollary}

\begin{proof}
 Let $\tilde{f}$ be a vector field satisfying~\eqref{eq:descend_vf} and let $x$ be any point of
 $\MA$. Note that $q$ is surjective, so $q^{-1}(x) \neq \emptyset$ and we can choose an arbitrary
 $\di \in \DI$ such that $q(\di) = x$. Take any smooth map $s : \MA \to \RR$ and compute
 $\tilde{s} = s\circ q$. By construction, $\tilde{s}$ satisfies~\eqref{eq:descend}, so
 \begin{displaymath}
  \lie_{\tilde{f}}(s\circ q)(\di) = 
   \lie_{\tilde{f}}(s\circ q)(\gamma\bscdot\di) \quad \text{for all $\gamma \in \Gamma$} \;.
 \end{displaymath}
 Direct application of the definition of the pushforward yields
 \begin{equation} \label{eq:descend_vf_2}
  q_*\tilde{f}(\di,t) = q_*\tilde{f}(\gamma\bscdot\di,t) \quad \text{for all $\gamma \in \Gamma$ and all $t \in \RR$} \;.
 \end{equation}
 Now we can use~\eqref{eq:descend_vf_1} to define $f$, that is, we set $f(x,t) = q_* \tilde{f}(\di,t)$.
 Note that, because of~\eqref{eq:descend_vf_2}, $f$ is well-defined, independent of the particular
 choice on $\di$. Smoothness can be established by computing $f$ using local coordinates.
\end{proof}

\end{document}
